\chapter{Pole Signage}
\label{ch:signage}

A hundred poles across a map are a hundred identical poles. The \textbf{Utility Pole Sign} is the
tag that tells them apart: which station feeds this one, which feeder it is on, who last climbed it.

\begin{iefork}
Not in stock Immersive Engineering, and not invented either. All thirteen kinds are signs that hang
on real poles --- the shapes, the colours and what each one means are the \emph{Los Angeles
Department of Water and Power}'s and \emph{Southern California Edison}'s. They are told apart the
way the real ones are: by shape and colour, long before you are close enough to read the number.
\end{iefork}

One item, one recipe, thirteen plates. Which plate a sign is lives on the sign, not on the item, so
there is nothing to hunt for in a creative tab and nothing to craft twice.

\section{Hanging one}

Place it against the \textbf{side} of whatever holds the wires up --- a pole, a crossarm, a wall.
The plate bolts flat to that face and reads from the outside. There is no top or bottom placement: a
tag lying on a floor is not a thing that exists, and a sign whose support is broken comes down with
it, the way a torch does.

Signs have no collision box. Walking into a pole covered in them is walking into the pole.

\section{Choosing the plate}

\textbf{Hit a sign with an Engineer's Hammer} and it steps to the next of the thirteen. Keep hitting
and it walks the whole set and comes back round --- one gesture, no menu, and it works from a ladder.

\textbf{Sneak and hit it with the hammer} and the plate opens for writing:

\begin{itemize}
\item The \textbf{arrows} pick the kind, the same set the hammer walks.
\item One \textbf{text box} per line that kind carries --- one, two or three, or none at all for the
      line crossing diamond, which is a symbol rather than a label.
\item The \textbf{preview} shows the real plate with the real lettering on it, laid out by exactly
      the arithmetic the world uses. What is in that window is what ends up on the pole.
\end{itemize}

Enter and Tab step between the lines. Done and Escape both keep what was typed.

\begin{ietip}
Text is \textbf{printed to fill the plate} rather than typeset at a fixed size. A short number comes
out big and a long one comes out small, which is what the real ones do and why a pole number and a
station name can share a plate the size of a hand. The vertical strips read downwards, because a
strip six pixels wide holds a pole number no other way.
\end{ietip}

A sign taken down keeps its plate and its text, so hammering thirteen times to find the one you
meant and then mining it by accident costs nothing.

\section{What the kinds mean}

\begin{description}
\item[Parallel Generation] The horizontal red strip. Several sources are feeding the distribution
      station this pole hangs off, or power is running in a loop across the area. Mainly a LADWP
      sign, and the only one in the set that carries white text. Two lines.
\item[Pole Number] The yellow vertical strip: general pole identification, and the one every
      utility uses. One line, read downwards.
\item[Street Light] The white vertical strip. SCE hangs these for street lighting, and so do
      privately owned poles. One line, read downwards.
\item[Old Pattern] The bare metal vertical strip --- mostly replaced by something more legible, and
      still on plenty of poles. One line, read downwards, in grey.
\item[Series Light Oval] The painted white oval a series-wired street light wears: the series number
      on top and the pole number underneath.
\item[Feeder Tag] The horizontal yellow strip. For the LADWP this is the distribution station, the
      feeder off it, and how many conduit or transformer bank connections are on that connection;
      on a light pole it is the capacitor bank or the phase cutoff. Everybody else uses it for pole
      numbers.
\item[Fibre Run] The orange strip, horizontal or vertical, marking fibre and cable runs. Almost
      every utility uses these.
\item[Pole Inspection Tag] The round bolt-on tag: who inspected the pole on top, the year
      underneath. It also does duty as the tag saying where a wooden pole was logged.
\item[Tower Number] The plain yellow diamond the LADWP numbers transmission towers with, roughly
      every fourth tower. No border --- the number is painted straight onto it.
\item[Line Crossing] The same diamond with a black outline and a black cross, marking where a line
      crosses another, or a span nobody wants to walk: a valley, a ravine. It carries no text.
\item[Tower Identifier] The tower run, in two forms. The vertical one reads downwards in three
      parts --- the initials of the plant the run starts at, the tower's number, then, under a rule
      printed on the plate, the initials of the station it ends at. The horizontal one is one line,
      and is what the DC towers of the Pacific Intertie carry.
\end{description}

\section{What a sign costs to look at}

Nothing you will notice. The plate is an ordinary block model baked into the chunk mesh like any
other --- one of fifty-two, thirteen plates times four facings --- so a pole line of them costs what
a pole line of blocks costs.

Only the \emph{lettering} is drawn per frame, and only within forty-eight blocks. Past that the
plate is a pixel or two across and the number was never legible anyway, so nothing is drawn at all.
A blank sign draws nothing at any distance.
